\bookchapter{الفصل 19: تشخيص الأعطال والإنقاذ الذاتي}

\chapterintro{الهاتف مثل الإنسان، يمرض أحياناً. بطء، انهيار التطبيقات، شاشة سوداء، عدم الشحن... معظم هذه المشاكل ليست مخيفة كما تبدو، ويمكنك تجربة حلها بنفسك قبل إرسالها للصيانة. هذا الفصل هو “دليل الطوارئ لهاتفك” — جرّب الطرق منخفضة التكلفة أولاً، وإذا لم تنجح فابحث عن طبيب محترف.}

\booksection{أولاً: بطء الهاتف وارتفاع حرارته — أكثر الأسئلة شيوعاً}

لنكن صادقين، أن يصبح الهاتف بطيئاً بعد حوالي عام من الاستخدام، هذا طبيعي جداً. السبب ليس بالضرورة عطلاً في الهاتف، بل في أغلب الأحيان “تراكم الغبار” — ملفات غير ضرورية، وعمليات خلفية، ومساحة تخزين تمتلئ تدريجياً.

\booksubsection{قم بالتشخيص أولاً، لا تتسرع في الإصلاح:}

\begin{enumerate}
  \item \textbf{تحقق من مساحة التخزين}: اذهب إلى الإعدادات التخزين لرؤية المساحة المتبقية. إذا كانت أقل من 10\%، سيصبح النظام بطيئاً جداً لأنه يكافح لإفساح مكان. من تجربتي: حافظ على 15\%-20\% على الأقل من المساحة المتبقية ليعمل الهاتف بسلاسة.
  \item \textbf{تحقق من ترتيب استهلاك البطارية}: اذهب إلى إعدادات البطارية لرؤية ترتيب الاستهلاك خلال 24 ساعة الماضية. إذا استحوذ تطبيق ما على النصيب الأكبر لكنك لم تستخدمه كثيراً، فهو يعمل في الخلفية. ألغِ تثبيته أو قيّد نشاطه الخلفي.
  \item \textbf{استرجع ما فعلته مؤخراً}: هل أصبح بطيئاً بعد ترقية النظام؟ (البطء في الأيام القليلة الأولى بعد تحديث كبير طبيعي، فالنظام يعيد بناء الفهارس في الخلفية). أم بعد تثبيت تطبيق جديد؟ (جرّب إلغاء تثبيته).
  \item \textbf{انظر إلى بيئة الاستخدام}: إذا تعرّض الهاتف للشمس المباشرة، أو كان في غلاف سميك أثناء الشحن ولعب الألعاب، فالحرارة حتمية. الحرارة العالية تجعل المعالج يخفّض تردده تلقائياً، فتشعر أن الهاتف أصبح بطيئاً فجأة. دعه يبرد ثم أعد المحاولة.
\end{enumerate}

\textbf{طرق العلاج (من البسيط إلى المعقد)}: أعد تشغيل الهاتف (يحل 60\% من المشاكل) نظّف مساحة التخزين (احذف التطبيقات والفيديوهات غير الضرورية) إذا كان تطبيق ما يسبب المشكلة فامسح ذاكرة التخزين المؤقت له أو ألغِ تثبيته وأعد تثبيته الملاذ الأخير: اعمل نسخة احتياطية للبيانات ثم أعد الضبط على المصنع.

\booksection{ثانياً: انهيار التطبيقات وعدم الاستجابة — التطبيق لم يعد يطيع}

تطبيق ما ينفتح ثم ينهار، أو يتجمد فجأة أثناء الاستخدام. عملية المعالجة مثل استقبال الطبيب، خطوة بخطوة:

\begin{enumerate}
  \item \textbf{أغلقه وأعد فتحه}: اسحب التطبيق من الخلفية، ثم أعد فتحه. معظم الأخطاء المؤقتة تُحل بهذه.
  \item \textbf{حدّث التطبيق}: اذهب إلى متجر التطبيقات لرؤية إن كان هناك تحديث. أحياناً يكون السبب أن إصدار التطبيق قديم جداً وغير متوافق مع النظام.
  \item \textbf{تحقق من الصلاحيات}: بعض التطبيقات تحتاج صلاحية معينة للعمل بشكل صحيح، لكنك لم تمنحها. مثلاً تطبيق الكاميرا إذا لم تُمنح صلاحية الكاميرا ينهار عند الفتح. اذهب إلى الإعدادات، ابحث عن ذلك التطبيق، وامنح جميع الصلاحيات اللازمة.
  \item \textbf{امسح الذاكرة المؤقتة}: اذهب إلى الإعدادات إدارة التطبيقات ابحث عن التطبيق التخزين مسح الذاكرة المؤقتة. انتبه، هو "مسح الذاكرة المؤقتة" وليس "مسح البيانات"! مسح البيانات سيحذف حالة تسجيل الدخول والإعدادات، وكأنك أعدت ضبط التطبيق على المصنع.
  \item \textbf{ألغِ التثبيت وأعد التثبيت}: الملاذ الأخير. قبل إلغاء التثبيت تأكد من عدم فقدان بيانات التطبيق — لـ WeChat انسخ سجل المحادثات احتياطياً أولاً، وللألعاب تأكد من ربط الحساب برقم الهاتف.
\end{enumerate}

\booksection{ثالثاً: عدم الشحن أو الشحن البطيء — قد لا يكون ذنب البطارية}

عدم الشحن مشكلة تسبب القلق بسهولة، لكن في كثير من الأحيان السبب في الملحقات وليس عطلاً في الهاتف.

\begin{itemize}
  \item \textbf{افحص منفذ الشحن أولاً}: استخدم عود أسنان أو أداة صغيرة غير موصلة، وأزل برفق الغبار والوبر من منفذ الشحن. واجهت هذا بنفسي — سقط خيط من الجيب وسدّ منفذ الشحن، ظننت أن المشكلة في التلامس، لكن بعد إخراجه عاد للعمل.
  \item \textbf{جرّب كابلاً ورأس شحن آخر}: الكابلات تتلف أكثر من رؤوس الشحن، خاصة تلك التي تُلفّ كثيراً. جرّب كابلاً ورأس شحن مؤكدين الصلاحية للاختبار المتقاطع، لتحديد أي الملحقات معطل بسرعة.
  \item \textbf{مشاكل الشحن اللاسلكي}: إذا كنت تستخدم الشحن اللاسلكي، تحقق من سماكة غلاف الهاتف، أو وجود رقائق معدنية خلفه (مثل حامل مغناطيسي). وأيضاً، ضع الهاتف بدقة في منتصف لوح الشحن — إذا انحرف فقد لا يشحن.
  \item \textbf{الشحن البطيء}: تأكد من أن قدرة رأس الشحن كافية (iPhone للشحن السريع يحتاج على الأقل رأس شحن PD بقدرة 20 واط)، وأن الكابل يدعم الشحن السريع (ليس كل الكابلات تمرر أكثر من 3 أمبير)، وهل الهاتف ساخن جداً (الحرارة الشديدة تخفّض القدرة تلقائياً لحماية البطارية).
\end{itemize}

\booksection{رابعاً: عدم الاتصال بالشبكة، انقطاع Wi-Fi المتكرر}

خارج المنزل تعمل البيانات الخلوية بشكل جيد، لكن عند العودة لا تتصل بـ Wi-Fi؛ أو Wi-Fi متصل لكن لا توجد شبكة. افحص بهذا الترتيب:

\begin{enumerate}
  \item \textbf{تشغيل/إيقاف وضع الطيران}: اسحب مركز التحكم، شغّل وضع الطيران لمدة 5 ثوانٍ ثم أطفئه. هذا يجعل الهاتف يبحث عن الشبكة مرة أخرى، وهو فعال جداً.
  \item \textbf{"نسيان" Wi-Fi ثم إعادة الاتصال}: اذهب إلى إعدادات Wi-Fi، انقر علامة التعجب بجانب اسم الشبكة (في أندرويد الترس)، اختر "تجاهل هذه الشبكة" أو "نسيان"، ثم أعد إدخال كلمة المرور للاتصال.
  \item \textbf{أعد تشغيل الموجّه}: كثيرون لا يطفئون موجّه المنزل لشهور، ومع الوقت يصبح أقل استقراراً. افصل الطاقة لمدة 30 ثانية ثم أعد توصيلها، وتُحل معظم مشاكل Wi-Fi.
  \item \textbf{تحقق إن كانت المشكلة في الهاتف}: هل تتصل الأجهزة الأخرى بنفس Wi-Fi بشكل طبيعي؟ إذا كانت كلها طبيعية وهاتفك فقط لا يعمل، فقد يحتاج الهاتف لإعادة التشغيل أو إعادة ضبط إعدادات الشبكة (الإعدادات عام نقل أو إعادة ضبط إعادة ضبط إعدادات الشبكة، لن يحذف بياناتك، لكنه سيحذف جميع كلمات مرور Wi-Fi).
\end{enumerate}

\booksection{خامساً: مشاكل الشاشة: لمس غير حساس، وميض، آثار باقية}

\begin{itemize}
  \item \textbf{اللمس غير الحساس}: انزع الغلاف الزجاجي المقوّى أولاً لترى إن كان ذنبه — الغلاف الزجاجي الرخيص مع الوقت يؤثر فعلاً على اللمس. وجود ماء على الشاشة أو زيت على اليد يؤدي أيضاً إلى انحراف اللمس. إذا لم يكن أي منها، أعد التشغيل ثم حاول.
  \item \textbf{آثار/صور خافتة على الشاشة (احتراق الشاشة)}: عرض نفس الصورة لفترة طويلة (مثل إبقاء الشاشة مضاءة أثناء تصفح TikTok) يترك "أثراً" على شاشات OLED. هذه ليست عطلاً في الشاشة، بل مؤقت. أطفئ الشاشة لنصف ساعة وعادةً تتعافى.
  \item \textbf{وميض الشاشة}: إذا كان الوميض عند السطوع المنخفض، فقد يكون بسبب تعديل PWM، الشائع في شاشات OLED. اذهب إلى الإعدادات وابحث عن "تعديل DC" أو "منع الوميض"، وفعّله لتحسن كبير. ليست كل الهواتف لديها هذه الميزة.
  \item \textbf{منطقة على الشاشة سوداء/مضيئة تماماً}: هذا خطير، على الأرجح تلف مادي في الشاشة، يحتاج استبدال. الشاشات المستبدلة في مراكز غير رسمية، عادةً ألوانها وسطوعها أقل من الأصلية، وتؤثر على مقاومة الماء.
\end{itemize}

\booksection{سادساً: عدم الإقلاع، التعثّر عند شاشة الشعار — لا تذعر}

هذه الحالة تسرّع ضربات القلب، لكن تمهّل، جرّب هذه أولاً:

\begin{itemize}
  \item \textbf{اشحن لمدة نصف ساعة ثم حاول}: في كثير من الأحيان تكون البطارية قد نفدت تماماً. أوصل الشاحن الأصلي، انتظر 30 دقيقة على الأقل، ثم حاول التشغيل. بعض الهواتف بعد نفاد الطاقة تحتاج للوصول لمستوى معين لتستجيب.
  \item \textbf{إعادة التشغيل الإجباري}: لهواتف iPhone ذات الشاشة الكاملة: اضغط مرة زر رفع الصوت، مرة خفض الصوت، ثم اضغط مع الاستمرار الزر الجانبي للطاقة حتى يظهر شعار Apple (قد يستغرق 10-15 ثانية). الطريقة العامة لأندرويد: اضغط مع الاستمرار زر الطاقة لأكثر من 10 ثوانٍ حتى يهتز أو تضيء الشاشة.
  \item \textbf{التعثّر في واجهة الإقلاع وإعادة التشغيل المتكرر}: تسمى هذه الحالة "Apple الأبيض" (iPhone) أو "التعثّر عند الشعار" (أندرويد). إذا لم تنجح إعادة التشغيل الإجباري، قد تحتاج للدخول إلى وضع الاستعادة والإصلاح من الحاسوب. أوصل iPhone بالحاسوب وافتح iTunes/Finder، بعد الدخول لوضع الاستعادة اختر "تحديث" (وليس "استعادة")، سيعيد تثبيت النظام مع الاحتفاظ بالبيانات. أندرويد ادخل وضع Fastboot أو Recovery، اختصارات الأزرار تختلف بين العلامات، يُنصح بالبحث عن طراز هاتفك + "وضع الاستعادة".
  \item \textbf{إذا فشل كل ما سبق}: لا مفر من اللجوء إلى صيانة محترفة. إذا كنت قد نسخت بياناتك احتياطياً، فلا تقلق كثيراً.
\end{itemize}

\booksection{سابعاً: مشاكل الصوت: السماعة لا تعمل، صوت المكالمة منخفض}

\begin{itemize}
  \item \textbf{لا صوت من السماعة}: تأكد أولاً من عدم تفعيل وضع الصامت (المفتاح الجانبي في iPhone)، أو وضع عدم الإزعاج، أو خفض الصوت إلى الحد الأدنى. يبدو سخيفاً، لكن هذه الحالات تشكل 40\%. ثم تحقق من اتصال Bluetooth بسماعة أو مكبر صوت.
  \item \textbf{الطرف الآخر لا يسمعك في المكالمة}: تحقق من انسداد فتحة الميكروفون — الفتحة الصغيرة في أسفل الهاتف يسهل انسدادها بالغبار. نظّفها بفرشاة ناعمة بلطف. لا تستخدم إبرة أو عود أسنان للدفع بقوة، فقد تتلف الشبكة الواقية من الغبار.
  \item \textbf{صوت أجش/مشوّش}: إذا كان بسبب دخول الماء للسماعة — لا تستخدم مجفف الشعر بالهواء الساخن! امتص الماء السطحي بمنديل ورقي، ثم اتركه يجف طبيعياً. يُقال إن تطبيقات طرد الماء بالموجات الصوتية فعالة إلى حد ما (تشغيل صوت منخفض التردد لإخراج الماء بالاهتزاز)، يمكنك تجربتها.
\end{itemize}

\booksection{ثامناً: الصيانة الرسمية، تجنّب هذه الفخاخ}

إذا جربت كل ما سبق ولم ينفع، فلا بد من اللجوء إلى صيانة محترفة. لكن سوق الصيانة غامض، دعني أحدثك عن فخاخ وقعت فيها وتجاربي:

\begin{itemize}
  \item \textbf{ابحث عن الرسمي الحقيقي، لا "الذي يبدو رسمياً"}: اذهب إلى موقع العلامة التجارية (Apple.com.cn، Huawei.com.cn إلخ) وابحث عن قائمة "مزوّدي الخدمة المعتمدين". المتاجر غير المعتمدة التي تعلق شعار Apple في الشارع قد لا تستخدم ملحقات أصلية، وبعد الإصلاح تختفي مقاومة الماء فوراً.
  \item \textbf{ما يجب فعله قبل الإرسال للصيانة}: انسخ جميع البيانات احتياطياً (لا تعتمد على عدم فقدان البيانات أثناء الصيانة)، سجّل الخروج من جميع الحسابات المهمة (Apple ID، WeChat، Alipay)، أوقف "العثور على هاتفي". بعض مراكز الصيانة ستطلب كلمة مرور قفل الشاشة، وهذا خطر على الخصوصية — من الأفضل أن تسجّل الخروج من حساباتك بنفسك قبل الإرسال.
  \item \textbf{فخاخ استبدال الشاشة}: الشاشات غير الرسمية تنقسم إلى "شاشة مضغوطة أصلية" (زجاج أصلي مستبدل)، "شاشة مضغوطة لاحقاً" (زجاج غير أصلي)، "شاشة مجمّعة" (كلها غير أصلية). الأسعار تختلف بضعفين، لكن جودة العرض والمتانة تختلف كثيراً أيضاً. إذا كان الهاتف لا يزال في فترة الضمان، التزم بالطريق الرسمي.
  \item \textbf{فخاخ استبدال البطارية}: البطاريات غير الأصلية قد تبالغ في سعتها المعلنة، وبعد نصف سنة تتدهور أكثر من البطارية الأصلية بعد سنتين. سعر استبدال بطارية Apple الرسمي ليس أغلى بكثير من غير الرسمي، يُنصح بالطريق الرسمي.
  \item \textbf{فخاخ استرجاع البيانات}: بعد تلف الهاتف بالماء/السقوط، إذا قال لك مركز الصيانة "استرجاع البيانات يتطلب رسوماً إضافية"، احذر — بعض التجار غير الشرفاء يبالغون في حجم التلف لجمع رسوم أكثر. استرجاع البيانات الرسمي يقدّم السعر أولاً ثم العمل، وإذا لم ينجح لا تُحسب رسوم.
\end{itemize}

\begin{tcbwarn}{الحل الأول الشامل: أعد تشغيل الهاتف}
أياً كانت المشكلة التي تواجهها — بطء، انهيار تطبيقات، عدم اتصال بالشبكة، صوت غير طبيعي — أعد تشغيل الهاتف أولاً. يبدو هذا وكأنه تجاهل، لكن من تجربتي، على الأقل 60\% من المشاكل الصغيرة تُحل بإعادة التشغيل. الهاتف مثل الحاسوب، مع الاستخدام المطوّل تتراكم في الذاكرة أشياء فوضوية، وإعادة التشغيل تفرغها وتبدأ من جديد. لذلك، في المرة القادمة التي يواجه فيها هاتفك مشكلة، لا تتسرع في البحث عن “ماذا أفعل”، أعد التشغيل ثم انظر.

\end{tcbwarn}

\begin{tcbpractice}{احفظ الآن: هذه العمليات الإنقاذية يجب أن تعرفها جيداً}
\begin{enumerate}
  \item \textbf{إعادة التشغيل الإجباري}: لهواتف iPhone ذات الشاشة الكاملة: زر رفع الصوت زر خفض الصوت اضغط مع الاستمرار زر الطاقة حتى يظهر الشعار. أندرويد: اضغط مع الاستمرار زر الطاقة لمدة 10-15 ثانية حتى يهتز.
  \item \textbf{تنظيف منفذ الشحن}: ابحث عن أداة صغيرة (عود أسنان أو أداة إخراج شريحة SIM)، ونظّف برفق غبار منفذ الشحن.
  \item \textbf{إعادة اتصال Wi-Fi}: اذهب إلى إعدادات Wi-Fi و"انسَ" شبكتك الحالية، وأعد إدخال كلمة المرور للاتصال.
  \item \textbf{تشغيل/إيقاف وضع الطيران}: اسحب مركز التحكم، شغّل وضع الطيران، عُدّ 5 ثوانٍ في ذهنك، أطفئه، وانظر إن عادت الشبكة.
  \item التقط لقطات شاشة لخطوات هذه العمليات واحفظها، أو اكتبها في الملاحظات. عند مواجهة مشكلة لا تذعر وتبحث، بل افتح الملاحظات مباشرة.
\end{enumerate}

\end{tcbpractice}
